\chapter{Implementacja funkcji zarządzania projektami}
\label{chap:implementacja-funkcji}

Podstawę działania systemu PlanWise stanowią funkcje umożliwiające organizację i monitorowanie pracy projektowej. Obejmują one obsługę kont, tworzenie projektów, zarządzanie zespołem, definiowanie zadań i sprintów oraz prezentowanie postępu prac. Dane powstające podczas wykonywania tych operacji są następnie wykorzystywane przez mechanizmy wspomagania decyzji.

W niniejszym rozdziale przedstawiono implementację funkcji użytkowych, ze szczególnym uwzględnieniem współpracy interfejsu Angular z częścią serwerową. Omówiono przebieg najważniejszych operacji, reguły zmiany stanu i sposób prezentacji danych. Szczegóły algorytmów oceny ryzyka, priorytetyzacji, przydzielania zadań i estymacji kosztów przedstawiono w kolejnym rozdziale.

\section{Obsługa kont użytkowników i sesji}

Dostęp do części projektowej aplikacji wymaga uwierzytelnienia. Interfejs udostępnia widoki rejestracji, logowania i resetowania hasła, a stan sesji przechowuje magazyn \texttt{AuthStore}, z którego token dostępu jest dołączany do kolejnych żądań. Przy ponownym uruchomieniu aplikacji podejmowana jest próba odtworzenia sesji przez operację odświeżenia; jej niepowodzenie powoduje wyczyszczenie stanu lokalnego.

Dwie decyzje implementacyjne mają znaczenie wykraczające poza sam przebieg logowania. Adres e-mail jest przed zapisem normalizowany przez usunięcie skrajnych białych znaków i zamianę na małe litery, co czyni regułę unikalności konta jednoznaczną. Wylogowanie obejmuje natomiast zarówno unieważnienie sesji po stronie serwera, jak i usunięcie danych z pamięci klienta, ponieważ samo ukrycie widoków projektowych nie unieważnia tokenu odświeżania.

Uwierzytelnienie stanowi wspólny punkt wejścia do funkcji projektowych, lecz nie zastępuje kontroli dostępu do konkretnego projektu, wykonywanej osobno przy obsłudze operacji na jego danych.

\section{Tworzenie i zarządzanie projektami}

Projekt stanowi nadrzędny kontekst organizacji pracy. Zawiera nazwę, prefiks kluczy zadań, oznaczenie procesu, opcjonalną nazwę klienta, status oraz identyfikator właściciela. Podczas tworzenia sprawdzana jest dostępność prefiksu, wykorzystywanego następnie do budowania czytelnych oznaczeń zadań, a właściciel jest automatycznie dodawany do listy członków.

Wartości początkowe nadawane właścicielowi, obejmujące pojemność równą jeden i stawkę godzinową równą zeru, wymagają dostosowania, jeżeli dane zespołu mają być wykorzystywane w analizach kosztowych. Model projektu przewiduje ponadto stany \texttt{Active}, \texttt{OnHold}, \texttt{Completed} i \texttt{Archived}, przy czym archiwizacja zmienia wyłącznie status i nie blokuje samoistnie operacji na zadaniach.

Identyfikator wybranego projektu znajduje się w adresie i określa zakres danych pobieranych przez widoki backlogu, tablicy, zespołu i analiz. Pozwala to zachować jednoznaczny kontekst podczas nawigowania pomiędzy funkcjami, bez ponownego wybierania przedsięwzięcia.

\section{Zarządzanie zespołem i parametrami jego członków}

Członkostwo reprezentuje encja \texttt{ProjectMember}, zawierająca identyfikator projektu, opcjonalny identyfikator konta, adres e-mail, rolę, pojemność, stawkę godzinową oraz listę kompetencji. Opcjonalność identyfikatora konta pozwala zapisać członka zespołu na podstawie samego adresu, zanim zostanie on powiązany z użytkownikiem systemu; nie oznacza to jednak wysłania zaproszenia, ponieważ dodanie wpisu i dostarczenie wiadomości są odrębnymi operacjami.

Model ten wprowadza dwa ograniczenia istotne dla późniejszych analiz. Wpis członkostwa posiada pojedyncze pole roli, wobec czego nie odwzorowuje osoby pełniącej w projekcie kilka ról zawodowych; rozbudowa wymagałaby określenia, jak przypisywać wówczas stawki i dostępność. Pojemność wyraża natomiast udział dostępności, a nie liczbę punktów zadaniowych, przez co widok zespołu musi odróżniać wielkość przypisanych prac od deklarowanego zaangażowania.

Parametry te są wykorzystywane poza modułem: stawki i role zasilają estymację kosztów, a kompetencje i pojemność wpływają na propozycje przydziałów. Ich kompletność przekłada się zatem bezpośrednio na użyteczność rekomendacji.

\section{Tworzenie i edycja zadań}

Podstawową jednostką pracy jest encja \texttt{ProjectTask}. Zadanie jest przypisane do projektu i posiada tytuł, opis, status, priorytet, oszacowanie punktowe, opcjonalnego wykonawcę i termin. Model zawiera również wartość biznesową, przynależność do sprintu oraz parametr określający kolejność.

Podczas tworzenia zadania backend najpierw sprawdza dostęp użytkownika do projektu. Następnie pobiera prefiks projektu i kolejny numer zadania. Na tej podstawie budowany jest klucz, na przykład \texttt{PW-12}. Identyfikator techniczny pozostaje niezależny od tego oznaczenia.

Nowe zadanie otrzymuje stan \texttt{Backlog}. Jego pozycja jest wyznaczana na podstawie największej dotychczasowej wartości porządkującej w backlogu, powiększonej o stały odstęp. Dzięki temu zadanie może zostać dodane na końcu listy bez zmiany pozycji wszystkich pozostałych elementów.

W systemie rozróżniono trzy właściwości związane z ważnością i kolejnością pracy:

\begin{itemize}
    \item \texttt{Priority} --- kategoria priorytetu, przyjmująca wartości \texttt{Low}, \texttt{Medium}, \texttt{High} lub \texttt{Urgent};
    \item \texttt{BusinessValue} --- wartość biznesowa wykorzystywana w analizach;
    \item \texttt{Rank} --- parametr porządkujący zadania w widoku.
\end{itemize}

Właściwości te pełnią różne funkcje. Zmiana kategorii priorytetu nie jest tym samym co przesunięcie zadania w backlogu, a wartość biznesowa jest jednym z kryteriów rekomendowania kolejności.

Edycja zadania wykorzystuje operację częściowej aktualizacji. Dla wybranych pól zastosowano typ \texttt{Optional<T>}, który pozwala odróżnić brak zmiany od jawnego usunięcia wartości. Umożliwia to na przykład wyczyszczenie wykonawcy, terminu, oszacowania lub przypisania do sprintu.

Po stronie klienta szczegóły zadania są prezentowane w panelu otwieranym z tablicy. Zapis zmian powoduje aktualizację wybranego zadania i ponowne pobranie danych widoku. Użytkownik może dzięki temu edytować element bez opuszczania kontekstu bieżącej pracy.

\section{Podzadania, etykiety i zależności}

Zadanie może zawierać podzadania, posiadające własny tytuł i znacznik wykonania, oraz etykiety pozwalające klasyfikować pracę niezależnie od statusu i priorytetu. Podzadania nie zastępują pełnych zadań projektowych: stanowią element zadania nadrzędnego, a ich wykonanie opisuje postęp jego wewnętrznych czynności.

Zależności reprezentują obiekty \texttt{TaskLink}, zawierające identyfikator drugiego zadania oraz typ relacji. Podczas dodawania sprawdzane jest istnienie obu zadań i ich przynależność do tego samego projektu. Dane te zasilają harmonogramowanie, ocenę blokad i priorytetyzację, przy czym nie każde powiązanie informacyjne powinno być traktowane jako warunek rozpoczęcia pracy.

Sprawdzenie przynależności do projektu nie jest jednak uzupełnione kontrolą acykliczności grafu. Wykrywanie cykli pozostaje obszarem wymagającym rozbudowy i ma bezpośrednie znaczenie dla poprawności obliczeń harmonogramu omówionych w kolejnym rozdziale.

\section{Organizacja backlogu i cyklu życia sprintu}

Sprint zawiera nazwę, opcjonalny cel, daty graniczne i stan, przechodząc kolejno przez wartości \texttt{Planned}, \texttt{Active} i \texttt{Completed}. Dozwolone przejścia kontrolują metody domenowe, a próba wykonania operacji w niewłaściwym stanie zwraca błąd. Procedura rozpoczynania sprintu sprawdza dodatkowo, czy w projekcie nie istnieje już sprint aktywny; reguła ta jest realizowana na poziomie przypadku użycia, wobec czego jej odporność na żądania równoczesne wymaga osobnej oceny.

Powiązanie przynależności do sprintu ze statusem zadania zaprojektowano asymetrycznie. Zadanie przypisane do sprintu z backlogu otrzymuje stan \texttt{Todo}, natomiast usunięcie przypisania przywraca stan \texttt{Backlog} wyłącznie wtedy, gdy zadanie pozostawało w stanie \texttt{Todo}. Prace rozpoczęte i zakończone zachowują swój stan, dzięki czemu zmiana organizacji pracy nie powoduje utraty informacji o postępie.

Zakończenie sprintu aktualizuje jego stan i udostępnia podsumowanie przypisanych zadań, lecz nie przenosi automatycznie elementów niezakończonych do kolejnej iteracji. Dalsze ich zaplanowanie wymaga osobnej zmiany przypisania.

\section{Implementacja tablicy Kanban}

Tablica przedstawia zadania w trzech kolumnach odpowiadających stanom \texttt{Todo}, \texttt{InProgress} i \texttt{Done}. Backend zwraca zawartość obejmującą cały projekt, a frontend zawęża prezentację do wybranego sprintu, wszystkich sprintów albo zadań nieprzypisanych, przyjmując domyślnie sprint aktywny. Interfejs obsługuje ponadto grupowanie według statusu, umożliwiające zmianę stanu przez przeciąganie kart, oraz według wykonawcy, służące do przeglądania rozdziału pracy.

Przeciąganie zaimplementowano jako aktualizację optymistyczną: klient aktualizuje widok lokalny, wysyła żądanie z docelowym statusem i indeksem, a w razie błędu pobiera dane ponownie. Rozwiązanie to poprawia bezpośredniość reakcji aplikacji, lecz wymaga przywrócenia zgodności z serwerem, gdy operacja się nie powiedzie.

Po stronie backendu pozycja jest zapisywana w polu \texttt{Rank}. Przy umieszczeniu zadania pomiędzy dwoma elementami wyznaczana jest średnia ich wartości:

\begin{equation}
    r_{\mathrm{new}} =
    \frac{r_{\mathrm{prev}} + r_{\mathrm{next}}}{2}.
    \label{eq:rank-zadania}
\end{equation}

Na początku lub końcu kolumny stosowany jest stały odstęp. Pozwala to uniknąć każdorazowej numeracji wszystkich zadań. Wielokrotne wstawianie pomiędzy te same elementy jest jednak ograniczone precyzją zapisu liczby i może wymagać okresowego uporządkowania wartości.

Przy przejściu do stanu \texttt{Done} zapisywany jest czas zakończenia. Jeżeli zadanie zostaje ponownie przeniesione do innego stanu, pole to jest czyszczone. Informacja jest wykorzystywana przy tworzeniu wykresu spalania.

Kontrakt tablicy przewiduje miejsce na limit WIP, lecz backend zwraca obecnie wartość pustą. Implementacja zapewnia zatem wizualizację przepływu i zmianę statusów, ale nie egzekwuje konfigurowalnych limitów pracy w toku.

\section{Harmonogram, wykres Gantta i kamienie milowe}

Widok harmonogramu przedstawia zadania na osi czasu. Dane wiersza obejmują identyfikator i nazwę zadania, daty graniczne, zapas czasu, oznaczenie krytyczności oraz identyfikatory poprzedników.

Ręcznie wskazane terminy przechowywane są w module \texttt{Scheduling} jako odrębne elementy harmonogramu, co oddziela ustawienia planu od danych opisujących samo zadanie i pozwala uwzględniać je przy kolejnych obliczeniach jako nadpisania. Przed zapisaniem zmiany system przelicza bieżący harmonogram i sprawdza, czy proponowane rozpoczęcie nie narusza zależności od poprzedników; wykryte naruszenie kończy operację błędem. Udostępniono również walidację całego zestawu proponowanych przesunięć, przy czym kontrola ta obejmuje wyłącznie zależności, nie zaś kalendarze pracowników ani konflikty zasobowe.

Kamień milowy reprezentuje nazwany punkt kontrolny z określonym terminem. Nie posiada parametrów pracochłonności ani wykonawcy, służy zatem do oznaczania istotnych dat, a nie do reprezentowania kolejnego zadania roboczego.

Sposób wyznaczania terminów, zapasu czasu i krytyczności oraz ograniczenia przyjętego modelu omówiono w rozdziale dotyczącym mechanizmów wspomagania decyzji.

\section{Pulpit projektu i monitorowanie postępu}

Pulpit projektu agreguje dane potrzebne do wstępnej oceny stanu prac. Zapytanie \texttt{GetOverviewQuery} przygotowuje liczby zadań w poszczególnych stanach, sumę punktów, liczbę punktów zakończonych, informacje o aktywnym sprincie i listę zadań wymagających uwagi.

Do ostatniej grupy zaliczane są niezakończone zadania z terminem wcześniejszym od bieżącej daty. Jest to rozpoznanie przekroczenia terminu na podstawie danych operacyjnych. Nie stanowi prognozy przyszłego opóźnienia.

Zestawienie obciążenia łączy członków zespołu z przypisanymi zadaniami. Może dotyczyć całego projektu lub wybranego sprintu. Zwraca między innymi pojemność członka i sumę punktów przypisanych prac.

W analizowanej procedurze suma obejmuje zadania niezależnie od ich statusu. Wynik opisuje zatem wielkość przypisanego zakresu, a nie wyłącznie pracę pozostałą do wykonania. Rozróżnienie to jest istotne podczas interpretowania zakończonego sprintu lub projektu.

Osobnym elementem monitorowania jest wykres spalania sprintu. Implementacja tworzy linię idealnego spadku pozostałych punktów i serię wyznaczaną z zapisanych dat zakończenia zadań.

Dla bieżącego zbioru zadań sprintu pozostały zakres w dniu \(d\) można zapisać jako:

\begin{equation}
    R(d) =
    P -
    \sum_{\substack{i \in S_{\mathrm{done}}\\c_i \leq d}} p_i,
    \label{eq:burndown-implementacja}
\end{equation}

gdzie \(P\) oznacza sumę punktów zadań obecnie przypisanych do sprintu, \(S_{\mathrm{done}}\) zbiór zadań obecnie zakończonych, \(c_i\) datę zakończenia, a \(p_i\) oszacowanie punktowe. Brak oszacowania jest w tej agregacji traktowany jako zero.

Dla dni przyszłych seria rzeczywista nie jest uzupełniana. W przypadku zakończonego zadania bez zapisanej daty zakończenia implementacja przyjmuje datę rozpoczęcia sprintu jako wartość zastępczą.

Wykres ma ograniczenia wynikające z modelu danych. Nie korzysta z pełnej historii zmian zakresu i statusów. Ponowne otwarcie zadania usuwa jego bieżącą datę zakończenia, a zmiana przypisania lub oszacowania może wpłynąć również na wcześniej prezentowane wartości. Wykres należy więc interpretować jako rekonstrukcję na podstawie dostępnego stanu i dat zakończenia.

Uzupełnieniem wykresu spalania jest zestawienie prędkości zespołu, obejmujące sprinty zakończone, uporządkowane według daty zakończenia. Dla każdego z nich zwracana jest suma punktów zadań do niego przypisanych oraz suma punktów zadań zakończonych.

Zestawienie to podlega temu samemu ograniczeniu co wykres spalania. Obie wielkości są wyznaczane z bieżącego przypisania zadań do sprintów, a nie z zakresu zapisanego w chwili jego rozpoczęcia. Przeniesienie zadania do innego sprintu po zakończeniu poprzedniego zmienia zatem również wartości prezentowane dla okresów wcześniejszych. Prędkość należy więc traktować jako wielkość rekonstruowaną, a nie zarejestrowaną w momencie zamknięcia sprintu.

\section{Komentarze, aktywność i powiadomienia}

Komentarze umożliwiają prowadzenie dyskusji w kontekście zadania; autor wpisu jest ustalany na podstawie sesji, a dodanie komentarza zgłasza zdarzenie domenowe. Rejestr aktywności przedstawia wybrane zmiany zachodzące w projekcie, pozwalając rozpoznać ostatnie działania bez otwierania kolejno wszystkich zadań. Nie należy go jednak utożsamiać z dziennikiem audytowym: pełny audyt wymagałby określenia zakresu rejestrowanych operacji, trwałości wpisów i informacji pozwalających odtworzyć zmiany wartości.

Powiadomienia stanowią odrębne dane powiązane z odbiorcą, zawierające typ, komunikat, opcjonalne odniesienie do projektu i odnośnik do miejsca w aplikacji. Użytkownik może je pobierać i oznaczać jako przeczytane.

Istotne jest, że moduł powiadomień zapisuje dane we własnym kontekście bazodanowym, wobec czego ich utworzenie nie należy do transakcji operacji z innego modułu. Podstawowa zmiana i powiązane z nią powiadomienie mogą zatem zakończyć się różnymi rezultatami.

\section{Aktualizacja widoków z wykorzystaniem SignalR}

Praca kilku osób nad projektem wymaga przekazywania informacji o zmianach pomiędzy otwartymi instancjami aplikacji. Wykorzystano do tego SignalR z grupowaniem połączeń według projektu, przy czym backend dodaje połączenie do grupy dopiero po zweryfikowaniu dostępu użytkownika do danego przedsięwzięcia.

Przyjęty przebieg synchronizacji jest celowo pośredni: po zapisaniu zmiany serwer rozsyła powiadomienie, a zainteresowane klienty pobierają aktualny stan z API. SignalR pełni więc wyłącznie funkcję sygnału do odświeżenia danych, natomiast źródłem ich aktualnej postaci pozostaje API. Subskrypcje są usuwane przy niszczeniu widoku, co ogranicza wielokrotne obsługiwanie tych samych komunikatów.

Mechanizm ten nie rozstrzyga samodzielnie konfliktów równoczesnej edycji. Odebranie informacji o zmianie i ochrona przed nadpisaniem cudzych danych są odrębnymi zagadnieniami, wymagającymi osobnej weryfikacji.

\section{Integracja widoków z operacjami analitycznymi}

Funkcje organizacyjne są połączone z analizami przez wspólny kontekst projektu, dzięki czemu przejście od backlogu lub harmonogramu do widoków ryzyka, priorytetów i kosztów nie wymaga ponownego wprowadzania danych. Operacje wykonywane w tle obsługiwane są przez identyfikator zlecenia: klient inicjuje analizę, śledzi jej stan i pobiera wynik po zakończeniu, o którym może zostać dodatkowo powiadomiony przez SignalR.

Kluczowe dla interpretacji wyników jest rozdzielenie prezentacji propozycji od jej zastosowania. Widok harmonogramu udostępnia panel propozycji przydziałów oraz panel wyjaśnienia zawierający cel obliczeń, założenia, nazwę mechanizmu i czas wygenerowania, a zatwierdzenie uruchamia odrębną operację zmieniającą dane. Opis oczekiwanej korzyści pozostaje przy tym elementem uzasadnienia rekomendacji, a nie dowodem poprawy uzyskanym w eksperymencie.

Błąd analizy jest prezentowany jako niepowodzenie konkretnej operacji i nie zmienia sam przez się danych backlogu ani wyniku wcześniej zakończonej estymacji.

\section{Konfiguracja i przygotowanie środowiska}

Uruchomienie wymaga skonfigurowania połączenia z PostgreSQL, parametrów uwierzytelniania oraz adresu API wykorzystywanego przez frontend; mechanizmy korzystające z modelu językowego wymagają dodatkowo konfiguracji usługi zewnętrznej. W środowisku deweloperskim aplikacja uruchamia migracje i udostępnia dokumentację OpenAPI, a repozytorium zawiera konfigurację Docker Compose wspierającą przygotowanie środowiska.

Przy przygotowywaniu danych do oceny systemu należy uzupełnić nie tylko zadania i sprinty, lecz także wartości biznesowe, zależności, kompetencje, pojemności i stawki. Brak tych informacji ogranicza możliwości analiz, nawet gdy podstawowe operacje tworzenia i edycji działają poprawnie.

Opisane funkcje tworzą warstwę danych operacyjnych dla mechanizmów inteligentnego wspomagania decyzji. Ich implementacja określa zarazem, jakie informacje są dostępne dla algorytmów i jakich danych historycznych nie można odtworzyć bez dalszej rozbudowy systemu.